% 摘要：采用人工撰写稿（data/2026zy.docx），数据已逐条核对（见 paper/README.md）
%
% 摘要必须**占满且仅占一页**。做法：本页单独放大行距，使内容刚好铺满版心。
% 行距系数 1.45 是按"小四字号 + 约 30 行"估算的：
%   · 若溢出到第 2 页 → 把下面的 1.45 调小（1.35 / 1.30 …）
%   · 若页面下留白 → 调大（1.55 / 1.65 …）
% 这个系数只作用于摘要页（在 \begingroup…\endgroup 内），不影响正文。
\begin{center}{\zihao{2}\heiti 基于能量分流与滚动随机优化的微网购电调控研究}\end{center}

\vspace{0.5em}
\begin{center}{\zihao{4}\heiti 摘\quad 要}\end{center}
\vspace{0.3em}
\begingroup\zihao{-4}\linespread{1.45}\selectfont

随着光伏发电技术的发展，集成光伏发电、储能设备与外部电网的非孤岛式微网在居民小区、产业园区等场景中逐渐得到普及。然而，光伏出力和负荷需求具有较强的不确定性，如何在保障负荷供电的同时降低购电成本，成为微网运行调控中的重要问题。本文以降低购电成本、保障小区负荷供电为目标，建立外部电网、光伏、储能与负荷协同运行的微网调控基础模型。针对确定性调度、预测不确定性、日内预报更新和实时电价波动四种条件，逐步引入随机情景、滚动决策与价格预测方法，构建递进式优化体系。

针对问题一，在全天电价、负荷及光伏预测功率已知的条件下，通过分别刻画外网供负荷、外网充电、光伏充电、储能充放电及弃光等功率流向，建立混合整数线性规划模型，并利用分支定界法对全天计划购电与储能调度策略进行求解，得到的最优方案全天购电量为 59482.70 kWh，购电费用为 35126.95 元，全天弃光量为 0，储能能够充分利用峰谷电价差和富余光伏实现“谷充峰放”。模型能量平衡和储能状态递推残差均达到 \(10^{-12}\) 数量级，验证了结果的数值可行性。

针对问题二，考虑到每天 0:00 制定购电计划时未来真实负荷与光伏尚未实现，而供电不足将产生 5 倍电价的紧急购电费用，本文将计划购电与实际执行分离，并利用历史预测残差构造随机情景，建立 7 日滚动两阶段随机优化模型。通过样本平均近似方法描述预测误差，同时利用多日滚动视野保留储能的跨日价值。对 2025 年 2 月至 12 月逐日求解，总购电费用为 1480.69 万元，累计紧急购电量为 26.9 万 kWh。进一步通过完美信息基准分解成本来源发现，模型相对理想信息条件增加的成本主要来自负荷和光伏预测不确定性，说明提高预测质量是降低日前计划风险的关键。

针对问题三，针对 0:00、6:00、12:00 和 18:00 四个时点能够获得最新光伏预报的特点，在问题二模型基础上建立日内多阶段滚动调整模型。每次新预报到达后，仅重新优化尚未执行的购电计划，并以当前真实储电量连接相邻决策阶段，同时根据调增 1.5 倍电价、调减承担 0.5 倍违约费用的规则统一计算调整成本。采用四个预报时点后，全年总费用为 1448.52 万元，紧急购电量降至 18.5 万 kWh，表明日内预报更新能够利用不断增加的信息修正日前计划，减少预测偏差造成的高价紧急购电。

针对问题四，在问题二日前随机计划和问题三日内滚动调整模型的基础上，进一步考虑计划制定时实时电价未知的情况，将负荷、光伏和电价共同视为随机因素，并引入电价预测及历史价格误差情景。日内调整时，根据已实现的真实电价更新后续时段的价格预测；针对可能出现的负价格，设置计划购电上限及调增、调减互斥约束，避免无效超量购电和数学套利。结果表明，引入日内预报和价格更新后，总费用由 1632.05 万元降至 1522.57 万元，降幅为 6.71\%，说明在前述模型基础上逐步增加信息更新机制，能够进一步优化购电时机与储能充放电策略。

\par\vspace{1.2\baselineskip}
\noindent\textbf{关键词：}微网调度；混合整数规划；滚动优化调整；样本平均近似预测 SAA
\endgroup

\newpage
